\documentclass[11pt,a4paper,openany]{report}
\usepackage[T1]{fontenc}
\usepackage[utf8]{inputenc}
\usepackage{lmodern}
\usepackage[margin=23mm]{geometry}
\usepackage{amsmath,amssymb,amsthm,mathtools}
\usepackage{longtable,booktabs,array,calc}
\newcounter{none}
\usepackage{graphicx,xcolor,tikz}
\usepackage{microtype}
\usepackage[unicode,breaklinks,colorlinks,linkcolor=blue!45!black,urlcolor=blue!45!black]{hyperref}
\urlstyle{same}
\providecommand{\tightlist}{\setlength{\itemsep}{0pt}\setlength{\parskip}{0pt}}
\providecommand{\passthrough}[1]{#1}
\providecommand{\NL}{\operatorname{NL}}
\providecommand{\colim}{\operatorname*{colim}}
\setlength{\parindent}{0pt}
\setlength{\parskip}{.4em}
\setlength{\emergencystretch}{3em}
\setcounter{tocdepth}{1}
\setcounter{secnumdepth}{0}
\renewcommand{\chaptername}{Lesson}
\DeclareUnicodeCharacter{03BB}{\ensuremath{\lambda}}
\DeclareUnicodeCharacter{0393}{\ensuremath{\Gamma}}
\DeclareUnicodeCharacter{2192}{\ensuremath{\rightarrow}}
\DeclareUnicodeCharacter{2011}{-}
\DeclareUnicodeCharacter{202F}{\,}
\DeclareUnicodeCharacter{2212}{\ensuremath{-}}
\DeclareUnicodeCharacter{03B1}{\ensuremath{\alpha}}
\DeclareUnicodeCharacter{03B2}{\ensuremath{\beta}}
\DeclareUnicodeCharacter{03BC}{\ensuremath{\mu}}
\DeclareUnicodeCharacter{03A6}{\ensuremath{\Phi}}
\DeclareUnicodeCharacter{2260}{\ensuremath{\ne}}
\DeclareUnicodeCharacter{00D7}{\ensuremath{\times}}
\DeclareUnicodeCharacter{2082}{\ensuremath{_2}}
\DeclareUnicodeCharacter{2228}{\ensuremath{\vee}}
\DeclareUnicodeCharacter{2032}{\ensuremath{^{\prime}}}
\DeclareUnicodeCharacter{220E}{\ensuremath{\square}}
\title{Tori, maximal tori and their conjugacy}
\author{GPT-6.1 Sol (OpenAI)\\Codex, Ultra setting}
\date{October 2026}
\begin{document}
\hypersetup{pageanchor=false}
\maketitle
\pagenumbering{roman}
\hypersetup{pageanchor=true}
\noindent Written and self-checked by the writing AI. No human or independent review is claimed.
Original contributions are dedicated under CC0 1.0. Cited works retain their own licences.
\par\medskip
\noindent Online course: \url{https://kokunoyumeto.github.io/open-math-courses/courses/AG-RG/}.
\begingroup\small\tableofcontents\endgroup
\clearpage
\pagenumbering{arabic}
\setcounter{chapter}{0}
\chapter{Tori, maximal tori and their conjugacy}\label{AG-RG-01-tori-maximal-tori-and-their-conjugacy}

\emph{Written by GPT-6.1 Sol (OpenAI) in Codex, Ultra setting, October 2026. Self-checked by the writing AI, GPT-6.1 Sol (OpenAI). Public domain (CC0).}

A diagonal matrix is easy to understand because its eigenspaces separate the action into one-dimensional pieces. A torus is the group-scheme version of this diagonal action. Over a field that is not algebraically closed, or over a scheme, the eigenspaces may appear only after changing the base. We therefore keep track of which conclusions hold on the original base and which hold locally for the étale topology.

We assume affine group schemes, smooth morphisms, faithfully flat descent, and the character description of groups of multiplicative type. Basic references are Conrad's \emph{Reductive group schemes}, SGA 3, and Milne's \emph{Algebraic Groups}. The first two treat arbitrary bases; the third provides complementary field arguments.

\section{Diagonal actions and their descent}\label{AG-RG-01-diagonal-actions-and-their-descent}

For a scheme \(S\), a \textbf{torus} is a group scheme \(T\) that is étale locally isomorphic to \(\mathbf G_{m,S}^{r}\), with \(r\) locally constant. Its character sheaf is

\[
X^*(T)=\underline{\operatorname{Hom}}_{S\text{-groups}}(T,\mathbf G_m).
\]

It is a locally constant sheaf of finite free abelian groups on the étale site. Its dual is the cocharacter sheaf \(X_*(T)\). The pairing is fixed by

\[
\chi\circ\lambda(t)=t^{\langle\chi,\lambda\rangle}.
\]

For \(T=\mathbf G_m^r\), the characters \(e_i(t_1,\ldots,t_r)=t_i\) and coordinate cocharacters form dual bases. A representation of a split torus on a finite locally free module \(M\) is a grading

\[
M=\bigoplus_{\chi\in\mathbb Z^r}M_\chi,
\]

with finitely many nonzero summands locally on \(S\). Indeed, write its comodule map as \(m\mapsto\sum_\chi m_\chi\otimes z^\chi\). Coassociativity implies that each \(m_\chi\) has weight \(\chi\), and the counit gives \(m=\sum m_\chi\). Linear independence of the Laurent monomials makes the sum direct. Each summand is a direct summand of \(M\), hence finite locally free.

\textbf{Example.} Multiplication on \(\mathbb C\) gives an embedding

\[
\operatorname{Res}_{\mathbb C/\mathbb R}\mathbf G_m\longrightarrow\operatorname{GL}_{2,\mathbb R},\qquad
a+bi\longmapsto\begin{pmatrix}a&-b\\b&a\end{pmatrix}.
\]

After base change to \(\mathbb C\), the two eigenvalues are \(a+bi\) and \(a-bi\), so this is a two-dimensional torus. Complex conjugation interchanges its two characters. A split real torus has trivial conjugation action on its character lattice, so this torus is not split. The descent action records a real distinction that disappears over \(\mathbb C\).

The anti-equivalence between tori and locally constant free character sheaves follows by descent from the split case: a change of basis in \(\mathbb Z^r\) gives the corresponding monomial automorphism of \(\mathbf G_m^r\), and these constructions carry inverse descent data to inverse descent data. The alternative definition using an fpqc splitting cover gives the same tori; we prove this after Lemma 3.1. One cannot replace the character sheaf by a single lattice with trivial action on an arbitrary base.

\section{1. What maximality means}\label{AG-RG-01-what-maximality-means}

A closed subtorus \(T\subset G\) is \textbf{maximal} if \(T_{\bar s}\) is a maximal torus of \(G_{\bar s}\) at every geometric point of \(S\). We always use this fibrewise definition. A \textbf{reductive group scheme} is a smooth affine group scheme of finite presentation whose geometric fibres are connected reductive algebraic groups. A connected smooth affine group over an algebraically closed field is reductive when it has no nontrivial smooth connected normal unipotent subgroup.

This definition of maximality is preserved by every base change. Over an algebraically closed field, maximality also persists on extending that algebraically closed field. To see this, spread a putative larger torus and its inclusion out over a finitely generated subalgebra of the extension field. The dimensions are locally constant after restricting to a suitable open subset. A closed point in that open subset has residue field the original algebraically closed field, and specialization would give a larger torus there.

\textbf{Example 1.1. Diagonal matrices.} In \(\operatorname{GL}_{n,S}\), the diagonal torus \(D=\mathbf G_m^n\) is maximal. For a matrix \(A=(a_{ij})\) over an \(S\)-algebra \(R\), commuting with the universal diagonal matrix says

\[
a_{ij}(z_i-z_j)=0\quad\text{in }R[z_1^{\pm1},\ldots,z_n^{\pm1}].
\]

Comparison of distinct Laurent monomials forces \(a_{ij}=0\) for \(i\ne j\). Thus \(C_{\operatorname{GL}_n}(D)=D\) as group schemes over every base. Any torus containing \(D\) centralizes it and is therefore equal to it.

\textbf{Example 1.2. Two real tori.} In \(\operatorname{SL}_{2,\mathbb R}\), compare

\[
D=\{\operatorname{diag}(t,t^{-1})\},\qquad
K=\left\{\begin{pmatrix}a&-b\\b&a\end{pmatrix}:a^2+b^2=1\right\}.
\]

Over \(\mathbb C\), the second group has characters \(a+bi\) and its inverse, so both groups become the same type of one-dimensional torus. An eigenbasis gives a conjugation in \(\operatorname{GL}_2(\mathbb C)\); rescaling one eigenvector makes its determinant one, giving a conjugation in \(\operatorname{SL}_2(\mathbb C)\). They are maximal because their complex fibres are maximal. They are not conjugate over \(\mathbb R\): the character lattice of \(D\) has trivial conjugation action, while that of \(K\) has the sign action. A real conjugation would be an isomorphism of real tori and would preserve these actions.

\section{2. Moving tori over an algebraically closed field}\label{AG-RG-01-moving-tori-over-an-algebraically-closed-field}

We first establish the classical tools used in the conjugacy argument. The proofs work in every characteristic. Orbit representability and the closed-orbit lemma are the exact supporting results of the supporting lesson \href{https://kokunoyumeto.github.io/open-math-courses/courses/AG-GS/AG-GS-03.html}{Group schemes over a field}; finite free-action quotients are provided by \href{https://kokunoyumeto.github.io/open-math-courses/courses/AG-GS/AG-GS-04.html}{Quotients and torsors}.

\textbf{A line detecting a subgroup.} Every closed subgroup \(H\) of an affine algebraic group \(K\) is the stabilizer of a line in a finite-dimensional representation. To construct it, let \(I\) be the defining ideal of \(H\) in \(k[K]\), with the right regular action. The coalgebra identities put any finite collection of functions in a finite-dimensional stable subspace: expand their coproducts in linearly independent first factors; coassociativity makes the second factors a stable finite-dimensional space. Choose such a space \(W\) containing ideal generators, and put \(V=I\cap W\). A group element stabilizes \(V\) exactly when its right translation preserves the ideal generated by \(V\), namely \(I\), which is exactly the condition that it belongs to \(H\). The line \(\bigwedge^{\dim V}V\) in \(\bigwedge^{\dim V}W\) has the same stabilizer. This is an identity on test algebras, since the ideal and exterior-power conditions are scheme-theoretic. Applying the same finite-dimensional construction to algebra generators of \(k[K]\) gives a faithful representation: its matrix coefficients generate \(k[K]\), so its map to a general linear group is a closed immersion. These observations prove the stabilizer and linearity statements used below.

\textbf{Lemma 2.A (all representations of a unitriangular subgroup).} Over any field \(k\), let \(N\) be a closed subgroup scheme of the upper unitriangular group over \(k\). Every finite-dimensional representation of \(N\) has a filtration whose nonzero successive quotients are one-dimensional trivial representations. Consequently \(N\) has no nontrivial character, and a diagonalizable subgroup scheme of \(N\) is trivial.

\textbf{Proof.} For the defining representation, the coordinate flag has trivial quotients. Tensor products inherit this property by ordering tensor basis vectors lexicographically. Direct sums inherit it by concatenating flags. A subrepresentation inherits the intersections with a flag; a quotient inherits its images, deleting zero quotients. These statements use ordinary vector-space intersections and remain statements about comodules, hence about every test algebra.

The coordinate algebra \(k[N]\) is generated by the upper matrix entries. Its right regular coaction sends a row by \[
\Delta(x_{ij})=\sum_l x_{il}\otimes x_{lj}.
\] The span of a row is therefore a quotient of the defining representation. Multiplication is a comodule map. The finite-dimensional span of polynomials of degree at most d in the entries is consequently a quotient of a finite direct sum of their tensor powers of orders at most d, and has the same trivial-quotient filtration.

Let \(W\) be an arbitrary finite-dimensional comodule. Its coaction embeds it in \(W_{\mathrm{triv}}\otimes k[N]\): the counit is a left inverse to the underlying linear map, and coassociativity makes it a comodule map when \(N\) acts only on the second factor of its target. Finitely many matrix coefficients of this map occur, and all belong to one of the preceding finite-degree spaces. Thus \(W\) embeds into a finite direct sum of such spaces and inherits the required filtration.

An invariant line in \(W\) has trivial action, by intersection with that filtration. Applied to the group-like function of a character in \(k[N]\), this proves character triviality as a coordinate-algebra identity, including nonreduced \(N\). If a diagonalizable subgroup \(D\) embeds in \(N\), each representation of \(D\) splits into weight spaces by AG-GS-05, Theorem 3.1. Its faithful defining representation has only trivial constituents by the filtration just proved, so it is trivial. Faithfulness gives \(D=1\). \(\square\)

A subgroup of multiplicative type cannot have a faithful unitriangular representation: its representations split into character spaces after the splitting cover, whereas the same filtration has only trivial quotients. Its action is therefore the identity. Consequently a subgroup of multiplicative type and a unitriangular subgroup of a linear group have trivial scheme-theoretic intersection. This conclusion includes finite nonreduced groups of multiplicative type.

\textbf{Triangularization.} Every finite-dimensional representation of a smooth connected solvable group \(B\) has an invariant full flag. Induct on the derived length. The connected reduced derived subgroup \(D\) has smaller derived length, so has common eigenspaces; the nonzero ones form a finite collection of linearly independent character spaces. Normality of \(D\) makes \(B\) permute them, and connectedness makes every permutation constant. Choose one, \(V_\chi\). The subgroup \(D\) acts on it by scalars. Its scalar character is trivial: each commutator has determinant one, so \(\chi^d=1\), where \(d=\dim V_\chi\); a smooth connected group has no nontrivial character with finite image. The last statement is valid when the finite image is nonreduced, because a map from a reduced scheme factors through its reduction. Thus the image of \(B\) on \(V_\chi\) is commutative. Commuting matrices have a common eigenvector: an eigenspace of a nonscalar member is a proper invariant space, and induction on its dimension finishes the argument. We have obtained an invariant line. Repeat on its quotient to obtain the full flag. This proves Lie--Kolchin triangularization directly, without using the fixed-point theorem.

\textbf{Fixed points on complete varieties.} A smooth connected solvable \(B\) acting on a nonempty complete variety has a fixed point. Choose a closed reduced orbit by \href{https://kokunoyumeto.github.io/open-math-courses/courses/AG-RG/AG-RG-S08.html}{Supporting group-scheme proofs}, Theorem G.4.2, whose hypotheses hold for this smooth solvable group and complete variety. Its stabilizer \(H\) is detected by the preceding line construction, so this complete homogeneous orbit maps equivariantly to projective space as the orbit of that line. Its image is closed. Triangularize the representation, with flag \(V_1\subset\cdots\subset V_n\). Choose the least \(d\) for which the closed image meets \(\mathbf P(V_d)\). The intersection avoids \(\mathbf P(V_{d-1})\), so is both complete and contained as a closed subvariety in an affine chart. It is therefore finite. Connectedness makes \(B\) fix its points. Since the image is one homogeneous orbit, that orbit is a point, and so was the original closed orbit. This proves the required fixed-point theorem.

\textbf{Solvable structure.} We prove that \(B=U\rtimes T\) for a smooth connected solvable \(B\) over our algebraically closed field, with \(U\) smooth connected unipotent and \(T\) a torus, and that all maximal tori are \(U\)-conjugate. Triangularize a faithful representation. Projection to its diagonal gives a torus \(D\) and a kernel \(U\) contained in the upper unitriangular group. Order the upper entries by increasing distance from the diagonal, refining equal distances in any order. Setting successive entries equal to zero gives a normal filtration of the unitriangular group with additive one-dimensional quotients: every cross term contributing to a given distance uses smaller distances, which have already vanished. Intersect it with \(U\). Each successive quotient is a closed subgroup of \(\mathbf G_a\), and diagonal conjugation on it is multiplication by a character of \(D\); the upper unitriangular part acts trivially. Remove repeated terms.

Here is the splitting step, including finite additive kernels. A proper closed subgroup \(N\) of \(\mathbf G_a\) is finite and is the kernel of an additive polynomial \(P\). Indeed its ideal has one monic generator of degree \(d\), with \(P(0)=0\). The subgroup identity makes \(P(X+Y)\) vanish modulo \((P(X),P(Y))\). The polynomial \(P(X+Y)-P(X)-P(Y)\) has degree less than \(d\) in each variable, so its vanishing in that quotient, whose basis is \(X^iY^j\) for \(i,j<d\), makes it zero. Comparison of coefficients gives only powers \(X^{p^i}\) in positive characteristic, or a linear polynomial in characteristic zero. The quotient is \(\mathbf G_a\) via \(P\). An action of a torus on this quotient is again linear, since automorphisms of the affine line fixing zero have degree one.

An extension of an affine torus \(D\) by a vector line has a regular scheme section. Here is a proof for its fppf vector-line torsor over \(\operatorname{Spec}A=D\). Choose a faithfully flat affine cover trivializing the torsor, with coordinate ring \(B\). On the double overlap the two line coordinates differ by translation by \(c\in B\otimes _A B\). The torsor relation makes c satisfy the additive cocycle identity on the triple overlap. Consequently the matrices \(\begin{pmatrix}1&c\\0&1\end{pmatrix}\) give descent data on \(B^2\). The affine module-descent proof in AG-GS-04, Section ``Faithfully flat affine descent'', or AG-RG-S04, Section A, descends them to an \(A\)-module \(E\) and an exact sequence \(0\to A\to E\to A\to0\): exactness is checked after the faithfully flat cover. Choose a preimage of \(1\in A\) under the surjection \(E\to A\). Its local coordinate vectors are (a,1), and the displayed transition matrices give exactly the torsor translation law. They therefore define a descended section of that torsor. This proof supplies the fppf assertion directly and needs no unproved comparison of fppf and Zariski cohomology. The multiplication defect is a regular degree-two cocycle. Exactness of torus invariants makes the regular group cohomology in positive degrees vanish, so a correction makes the section a homomorphism. The same vanishing in degree one makes two sections conjugate by the vector line. The regular cochain contraction is proved in AG-GS-05, Proposition 7.14, and the homomorphism-lifting and conjugacy calculation is AG-GS-05, Theorem 7.15.

An extension by a finite \(D\)-stable \(N\subset\mathbf G_a\) also splits. Push it out along \(N\hookrightarrow\mathbf G_a\); concretely this is the finite free-action quotient of the product with \(\mathbf G_a\). The resulting vector-line extension splits by the preceding paragraph. In these coordinates the original extension is described by a crossed homomorphism \(D\to\mathbf G_a/N\): in each fibre its points are a coset of \(N\). That crossed homomorphism is principal, by degree-one vanishing for the quotient vector line. Its conjugating element lifts from \((\mathbf G_a/N)(k)\) to \(\mathbf G_a(k)\), since \(k\) is algebraically closed. After conjugation the original extension is \(N\rtimes D\). This also covers \(N=\mathbf G_a\) by the preceding paragraph.

Induct on the filtration length of \(U\). In \(B/N\) a section of the diagonal quotient exists by induction; its inverse image is the extension by the last, central-in-\(U\), additive subgroup \(N\) just treated. Splitting that extension gives a section in \(B\). Quotients by the possibly nonfinite unipotent terms are affine: for a normal unipotent subgroup, Lemma 2.A makes the character of the detecting line on that subgroup trivial, and its translates span a representation on which the subgroup acts trivially. The representation kernel is exactly the subgroup; its closed image realizes the quotient. The quotient map is faithfully flat, by generic flatness and translation over the field. Thus the induction uses actual algebraic groups.

We now have \(B\simeq U\times D\) as varieties. Smoothness and connectedness of \(B\) and \(D\) imply smoothness and connectedness of \(U\). The torus section is \(T\).

\textbf{Lemma 2.B (the smooth central filtration).} Let \(B\) be a smooth connected solvable affine group with the splitting \(B=U\rtimes T\) just constructed, where \(U\) is smooth connected and unitriangular. There is a \(B\)-stable filtration \[
U=U_0\supset U_1\supset\cdots\supset U_a=1
\] by smooth connected normal subgroups such that every successive quotient is \(\mathbf G_a\), each \(T\)-action on it is multiplication by a character, and \[
[U,U_i]\subset U_{i+1}.
\]

\textbf{Proof.} Use the splitting \(B=U\rtimes T\) just proved. In the faithful triangular representation choose a \(T\)-weight basis compatible with the full \(B\)-stable flag: exactness of the \(T\)-weight projections lets one choose a weight basis in each flag quotient and lift it inside that flag term. In this basis \(T\) is diagonal and \(B\) remains upper triangular. In its unitriangular group, order the entries by increasing distance \(j-i\) from the diagonal and refine each distance arbitrarily. Let \(F_i\) be the subgroup on which the first i entries vanish. The multiplication formula for an entry of distance r has cross terms of smaller positive distances. On \(F_i\), all distances smaller than the distance currently removed have vanished. The next entry therefore gives a homomorphism \(F_i\to\mathbf G_a\) with kernel \(F_{i+1}\). Conjugation by a unitriangular matrix changes an entry only by terms involving smaller distances; hence each \(F_i\) is normal and \([\mathrm{UT},F_i]\subset F_{i+1}\). Diagonal conjugation scales the next entry by a character. Thus the entire upper triangular group preserves the filtration.

Put \(V_i=(U\cap F_i)^0_{\mathrm{red}}\). Over the perfect \(k\) these are smooth connected subgroups. They are \(B\)-stable: conjugation preserves \(U\cap F_i\), and the morphism from the smooth connected \(B\times V_i\) has image in its reduced identity component. Similarly, the commutator morphism \(U\times V_i\to U\cap F_{i+1}\) has smooth connected reduced source and contains the identity in its image, so factors through \(V_{i+1}\). This gives the central-filtration assertion before any claim about the quotients.

The quotient \(V_i/V_{i+1}\) exists, is smooth connected and affine by AG-GS-04, Theorem 11.1d. Smoothness can also be checked after algebraic closure: faithful flatness into its coordinate algebra makes the quotient geometrically reduced; a reduced algebraic group over the perfect field is smooth by AG-GS-02, Theorem 5.1. The reduced identity component of the kernel of the next-entry map on \(V_i\) is \(V_{i+1}\): the reduced identity component of \(U\cap F_{i+1}\) is contained in \(V_i\), and taking this intersection introduces no new reduced identity component. The quotient of that kernel by \(V_{i+1}\) is therefore zero-dimensional and finite, retaining its possible infinitesimal thickening. The image of the entry map has dimension at most one. The torsor dimension formula for the map onto its schematic image, AG-GS-03, Proposition 5.8 and Theorem 7.11, now gives \(\dim(V_i/V_{i+1})\le1\). If the dimension is zero, smooth connectedness makes the quotient trivial, so \(V_i=V_{i+1}\).

Otherwise the quotient is a one-dimensional smooth connected unipotent group. Here unipotence does not require an unproved quotient criterion: a faithful representation of the affine quotient, pulled back to \(V_i\), has a trivial-quotient filtration by Lemma 2.A; faithful flatness of \(V_i\) onto the quotient makes that filtration a quotient-group filtration as well. Its faithful image is unitriangular. The completed-curve and translation-family argument in earlier AG-RG-S03, Theorems 2.2 and 3.1 and Section 4, now identifies this quotient with \(\mathbf G_a\) as a group scheme.

The \(T\)-action is a regular action by automorphisms of this affine line fixing zero. On each closed point of \(T\) the automorphism is a linear polynomial. In the universal action its coefficients of higher powers therefore vanish on all closed points of the reduced torus, and are zero. Its slope is a unit and the action identity makes that slope a character of \(T\). Removing the repeated \(V_i\) gives the required filtration. This argument retains finite and infinitesimal kernels of the original entry maps; it does not mistake those maps for isomorphisms. \(\square\)

Here are the degree-one and degree-two calculations explicitly. For \(D=D_k(M)\), let \(\int:k[M]\to k\) extract the coefficient of the identity character. It is translation invariant on every scalar extension because \(\Delta(e^m)=e^m\otimes e^m\). If \(V\) is a vector representation and c is a regular cocycle \[
c(ab)=c(a)+a\,c(b),
\] put \(v=\int_b c(b)\). Integration of this equality in b gives \[
v=c(a)+a v,\qquad c(a)=v-a v.
\] This is an equality of regular maps, also when \(D\) is finite nonreduced. It divides by no group order. For a regular two-cocycle c, the identity \[
a c(b,d)-c(ab,d)+c(a,bd)-c(a,b)=0
\] and integration in d give \(c(a,b)=a h(b)-h(ab)+h(a)\), with \(h(a)=\int_d c(a,d)\). These are the two calculations used in the splitting argument.

\textbf{Lemma 2.C (unipotent cocycles on all test schemes).} Let a diagonalizable \(D\) act on a smooth connected \(U\) with a \(D\)-stable filtration of the kind in Lemma 2.B. Every regular crossed homomorphism \(D\to U\) is principal. More generally, over any \(k\)-algebra \(R\) it is principal after an fppf extension of \(R\).

\textbf{Proof.} Project the cocycle to \(U/U_1\), a vector line. The displayed degree-one identity makes its projection principal with parameter v. Over \(k\), the nonempty fibre of \(U\to U/U_1\) above v has a \(k\)-point: it is finite type and \(k\) is algebraically closed. Over \(R\), that fibre is a torsor under \(U_1\) and supplies an fppf covering on which v has a lift u. Replace the cocycle by \[
c'(d)=u^{-1}c(d)\,d(u).
\] Its projection to \(U/U_1\) is now one, so it takes values in \(U_1\) and still satisfies the crossed-homomorphism identity. Repeat through the finite filtration. Over \(k\) no extension is required. Over \(R\) the finitely many fppf refinements remain an fppf cover, and the product of the correcting lifts makes the original cocycle principal. All maps and equations used are comodule identities over the universal \(D_R\), rather than statements about \(D(k)\). \(\square\)

The filtration is central in \(U\), and each \(T\)-action on a quotient is linear. A torus \(M\subset B\) intersects \(U\) trivially, as its diagonalizable representation cannot have a nontrivial unitriangular image. Its projection embeds in \(D\). In \(U\rtimes q(M)\) it is a complement. Lemma 2.C conjugates this complement to the chosen section by an element of \(U\). Maximality then makes \(q(M)=D\). This proves both the structure theorem and conjugacy of maximal tori. In a commutative \(B\), the action is trivial, so the decomposition is the direct product \(T\times U\).

\textbf{Lemma 2.1.} If \(G\) is smooth connected affine over an algebraically closed field \(k\), its maximal smooth connected solvable subgroups are conjugate, and \(G/B\) is complete for each such subgroup \(B\).

\textbf{Proof.} Let \(B\) be a smooth connected solvable subgroup of largest dimension. Its line-stabilizer representation can be enlarged by a faithful representation. A \(B\)-stable full flag beginning with the detecting line then has stabilizer exactly \(B\): a flag stabilizer is contained in the detecting line's stabilizer, while \(B\) fixes every member of the flag. In a faithful representation every full-flag stabilizer has solvable reduced identity component, since it lies in an upper triangular group. Its dimension is consequently at most \(\dim B\). Thus the orbit of our flag has minimum possible dimension.

An orbit is locally closed by the earlier programme orbit theorem. For a smooth connected acting group it is irreducible. Its complement in its closure is a proper closed invariant subset of strictly smaller dimension. Every orbit in that complement has strictly smaller dimension than the original orbit. Our minimum-dimension orbit can therefore have no boundary. It is closed in the projective flag variety and hence proper, by \href{AG-RG-S01.html\#projective-properness}{AG-RG-S01, Lemma 7.A and Corollary 7.B}. Its scheme stabilizer is \(B\), so the orbit theorem identifies it with \(G/B\).

For any other maximal smooth connected solvable subgroup \(B'\), the fixed-point theorem proved above supplies a fixed point \(gB\) for the action of \(B'\) on \(G/B\). The stabilizer condition says \(g^{-1}B'g\subset B\). Maximality of \(B'\) forces equality, and conjugation transfers properness to \(G/B'\). This completes the proof using only results actually written in this lesson or in the earlier orbit lesson. For additional reading, compare the freely accessible \href{https://www.jmilne.org/math/CourseNotes/iAG200.pdf}{Milne author draft, v2.00}, Theorem 4.19, with the detecting-line construction proved above. \(\square\)

\textbf{Lemma 2.D (the no-torus criterion).} A smooth connected affine group \(H\) over an algebraically closed field is unipotent if it has no nontrivial torus. Consequently any such \(H\) which is not unipotent has a positive-dimensional maximal torus.

\textbf{Proof.} Choose a Borel \(B\). The solvable structure just proved gives \(B=U\rtimes T\); the hypothesis makes \(T\)=1, so \(B\) is unitriangular. Detect \(B\) as the stabilizer of a line \(kv\) in a finite-dimensional \(H\)-representation. Its line character is trivial by Lemma 2.A, so \(B\) fixes v as a vector, and its schematic vector stabilizer is still \(B\). The map \(h\mapsto hv\) descends through the full fppf quotient \(H\to H/B\) to a morphism \(H/B\to V\). Lemma 2.1 realizes \(H/B\) as a closed subvariety of the projective flag variety. It is integral: its underlying space is the surjective image of the irreducible smooth connected \(H\), and its charts are reduced because their functions inject under the faithfully flat pullback into functions on reduced open subschemes of \(H\). AG-RG-S01, Theorem 8.4 and its closed-immersion cohomology calculation make \(H^0(H/B,\mathcal O)\) a finite-dimensional \(k\)-domain. A finite-dimensional domain is a field, and the algebraically closed \(k\) makes that field \(k\). Thus every coordinate of the morphism to \(V\) is constant. Its value is v at the identity coset, so \(H\) fixes v. The vector-stabilizer identity forces \(H=B=U\), proving the assertion. The contrapositive and existence of maximal tori over the algebraically closed field give the final sentence. \(\square\)

\textbf{Theorem 2.2.} Any two maximal tori in a smooth connected affine group over an algebraically closed field are conjugate.

\textbf{Proof.} A torus is connected and solvable, so it lies in a maximal connected solvable subgroup; one obtains such a subgroup by increasing dimension until this is no longer possible. For maximal tori \(T_1,T_2\), choose \(B_1,B_2\) containing them. Lemma 2.1 gives \(gB_2g^{-1}=B_1\). Both \(T_1\) and \(gT_2g^{-1}\) are maximal tori of the connected solvable group \(B_1\). The solvable-group conjugacy theorem supplies \(u\in B_1(k)\) with \(u gT_2g^{-1}u^{-1}=T_1\). \(\square\)

No conclusion here asserts conjugacy under \(G(k)\) when \(k\) is not algebraically closed. Example 1.2 is already a counterexample.

\section{3. Deforming an embedding rather than an eigenvalue}\label{AG-RG-01-deforming-an-embedding-rather-than-an-eigenvalue}

We need a mechanism that transports a torus from one fibre to nearby fibres. It uses the entire torus action, which retains information even when differentiating its characters loses information in positive characteristic.

For tori \(T_1,T_2\subset G\), define

\[
\operatorname{Transp}_G(T_1,T_2)(R)
=\{g\in G(R):g(T_1)_R g^{-1}=(T_2)_R\}.
\]

The normalizer \(N_G(T)\) is the transporter from \(T\) to itself. The centralizer \(C_G(T)\) asks that the induced automorphism of \(T\) be the identity. Both conditions concern subgroup schemes after arbitrary base change, rather than just the rational points of a fibre.

\textbf{Lemma 3.1. Torus deformation.} Let \(G\) be smooth affine of finite presentation over \(S\).

\begin{enumerate}
\def\labelenumi{\arabic{enumi}.}
\tightlist
\item
  The normalizer and centralizer of a subtorus are closed subgroup schemes of finite presentation. Transporters between tori of the same rank are schemes of finite presentation.
\item
  These normalizers, centralizers, and transporters are smooth over \(S\).
\item
  A torus embedded in a fibre \(G_s\) extends, after an étale neighbourhood with unchanged residue field, to a torus embedded in \(G\).
\end{enumerate}

\textbf{Proof.} We explain the equations, the deformation obstruction, and then the passage from formal to étale neighbourhoods.

After splitting a torus, its finite subgroups \(T[n]\) are schematically dense, universally on the base. A Laurent polynomial has finite support; take \(n\) larger than all differences between exponent coordinates. Its distinct monomials remain distinct modulo \(n\), so a nonzero Laurent polynomial remains nonzero on \(T[n]\). This works over any ring. The closed conditions on every \(T[n]\) therefore detect both inclusion and equality of maps on \(T\).

For a finite locally free scheme \(Y\) and an affine finitely presented scheme \(X\), maps \(Y\to X\) are represented by an affine finitely presented scheme. On an affine chart choose a basis of \(\mathcal O(Y)\); the images of generators of \(\mathcal O(X)\) are coordinate vectors, and its finitely many relations become polynomial equations in those vectors. Factoring through a closed subscheme adds closed equations. Apply this construction to conjugation on \(T[n]\). The resulting intersections represent the centralizer and the inclusion transporter. Reduction to a noetherian model makes the defining ideals finitely generated. Descent of the finitely many equations and their base-change interpretation then gives finite presentation over the original base. For equal-rank tori, an inclusion is an equality if it is so on one fibre, after restricting to a neighbourhood; the equality transporter is the corresponding open part. For the normalizer the two inclusion conditions give equality directly.

Now let \(R\to R/I\) have \(I^2=0\). Two liftings of a torus homomorphism \(f_0:T_0\to G_0\) differ by a regular crossed homomorphism with values in

\[
M=\operatorname{Lie}(G_0)\otimes_{R/I}I,
\]

with the action through \(\operatorname{Ad}\circ f_0\). More explicitly, lifting a scheme map first gives a multiplication defect \(c(t,u)\), and associativity is exactly the degree-two cocycle equation. Correcting a lift by a function \(b(t)\) changes \(c\) by its coboundary. Once the defect vanishes, the difference between two homomorphisms is a degree-one cocycle. Conjugating by an element infinitesimally equal to one changes that cocycle by a degree-one coboundary.

Taking invariants for a torus is exact, by the weight decomposition. The regular group-cochain resolution computes its derived functors, so its cohomology in positive degrees is zero. One may see the contraction directly in the homogeneous bar resolution: integration in one torus variable means taking the coefficient of its trivial character. This normalized invariant projection commutes with translations, and the alternating face maps give \(dh+hd=1\) in positive degrees. Thus the degree-two defect can be removed, and any two liftings are infinitesimally conjugate. This is the obstruction argument of SGA 3, Exposé III, §2, in the special case in which all obstructions vanish. The complete regular cochain contraction is AG-GS-05, Proposition 7.14; its application to homomorphism liftings and nilpotent conjugacy is AG-GS-05, Theorem 7.15.

Lift a transporter point \(g_0\) arbitrarily using smoothness of \(G\). The transported torus and the target torus have the same reduction. The degree-one calculation corrects \(g\) by an element reducing to one so that it transports the whole torus. For a centralizer point, the corrected element first normalizes \(T\). Its action on the character lattice reduces to the identity. Automorphisms of a torus are étale locally integral matrices, with no infinitesimal deformations, so this action is the identity already over \(R\). This proves formal smoothness of the three representing schemes, hence smoothness by finite presentation. In particular,

\[
\operatorname{Lie}(C_G(T))=\operatorname{Lie}(G)^T.
\]

Work on an affine base \(\operatorname{Spec}R\) at the chosen prime \(p\). Descend the finitely presented smooth \(G\) to a finitely generated \(\mathbf Z\)-subalgebra \(R_0\) using AG-MO-03, Theorems 4.1--4.2, and retain smoothness by the finite conormal splitting witnesses of AG-RG-S02, Lemma 3.1 and Proposition 3.2. Include the specified fibre torus in the same finite data before taking a local model. Its Hopf presentation, finite separable splitting algebra, splitting maps and inverses, descent maps, embedded coordinate pullback and finitely many surjectivity witnesses use only finitely many coefficients in \(\kappa(p)=\operatorname{Frac}(R/p)\). Express those coefficients as fractions of elements of \(R\) modulo \(p\) and include all their numerators and denominators in \(R_0\). With \(p_0=p\cap R_0\), the residue field \(\kappa(p_0)\) embeds in \(\kappa(p)\). Enlarge \(R_0\) until it contains these finite presentations, expressions and relation witnesses. Their identities then hold over \(\kappa(p_0)\), because this field embeds in \(\kappa(p)\); their chosen nonzero determinants remain nonzero. The resulting source is a torus split by the retained finite étale algebra and is embedded in the fibre of the smooth stage. Now take the Noetherian local ring \((R_0)_{p_0}\), essentially of finite type over \(\mathbf Z\), and its completion. The argument below lifts that exact stage fibre torus. After the final pointed étale approximation and pullback to \(R\), its residue data recover the originally specified torus and embedding, with unchanged residue field at \(p\).

Over its complete local ring \(A\), lift the fibre torus as follows. AG-RG-S05, Lemma 3.A constructs, from a finite Galois residue splitting field \(L/\kappa\), a finite free etale faithfully flat local algebra \(E/A\) with lifted Galois automorphisms and \(E\otimes_AE\simeq\prod_\Gamma E\). Descend \(E[M]\) with its fixed semilinear character-lattice action. That lemma proves that this is an actual finitely presented torus \(T_A\), and that every reduction \((T_A)_{A/\mathfrak m^n}\) is the descent of the same character data on \((E/\mathfrak m^nE)[M]\). Thus the source is already defined over \(A\). Apply AG-GS-05, Theorem 7.15 to this fixed source and the smooth target to lift the specified fibre homomorphism successively modulo \(\mathfrak m^n\), retaining each previous lifting. Its existence is supplied by the regular cochain contraction of AG-GS-05, Proposition 7.14. We now prove effectivity of the compatible homomorphisms.

Let \(A\) be Noetherian and complete for an ideal \(J\), and first let the source be \(D_A(M)\). Its completed coordinate module consists of families \((a_m)_{m\in M}\) tending \(J\)-adically to zero, while its actual coordinate module consists of families of finite support. Compatible homomorphisms modulo \(J^n\) give a map from \(A[G]\) to the former. For \(f\in A[G]\), write its coproduct as a finite sum \(\sum_{i=1}^r f_i\otimes g_i\). Compatibility with the coproduct says that the coefficient matrix of its completed image has entries

\[
\delta_{m,n}a_m=\sum_{i=1}^r b_{i,m}c_{i,n},
\tag{3.1}
\]

where the right side consists of the coefficient families of the images of the \(f_i\) and \(g_i\). Regard these columns as elements of \(A^M\). Every \(a_ne_n\) belongs to the finite module spanned by the \(r\) families \((b_{i,m})_m\). The submodule generated by all \(a_ne_n\) is therefore finitely generated, since \(A\) is Noetherian. It is the direct sum of the ideals \(Aa_n\) in their distinct coordinate positions. Finite generation forces all but finitely many of these ideals to be zero. Thus the image of \(f\) has finite support. This holds for every \(f\), so the completed map takes values in \(A[M]\) itself. Its algebra and Hopf identities hold because \(A[M]\) and \(A[M\times M]\) inject into their completions. This proves existence and uniqueness of the actual homomorphism. For the torus source just constructed, use the particular finite free cover \(E/A\) of AG-RG-S05, Lemma 3.A. Both \(E\) and \(E\otimes_AE\simeq\prod_\Gamma E\) are complete for the extended ideal, by their finite free bases; they are Noetherian because they are finite over the Noetherian ring \(A\). Apply the split argument over \(E\). Its two pullbacks to \(E\otimes_AE\) have the same reductions at every order, since the original formal maps respected the character descent data in (3.5) of that lemma. Uniqueness of the split effectivity argument makes those pullbacks equal as actual Hopf maps. AG-RG-S04, Lemmas A1.1--A1.2 therefore descend the actual homomorphism to \(A\), and faithful flatness verifies all of its Hopf identities and the prescribed reductions. This proves the effectivity statement needed here.

The lift is an embedding. Its kernel meets each finite subgroup \(T[n]\) in a finite scheme whose special fibre is trivial. The augmentation ideal of that finite algebra is a finite module over the complete local ring, and its reduction vanishes; Nakayama makes it zero. Hence the kernel is trivial on every \(T[n]\). On each geometric fibre the kernel is a diagonalizable subgroup of a torus, so density of its torsion implies that the fibre kernel is trivial. The identity section into the kernel is a finitely presented closed immersion with flat source and is an isomorphism on every fibre. It is an isomorphism: for its ideal \(I\) the exact sequence and flatness of the quotient give \(I\otimes k(s)\hookrightarrow A_K\otimes k(s)\); fibre equality makes this zero, and Nakayama locally gives \(I=0\).

Here is also the closed-immersion criterion used for this monomorphism. Split its multiplicative-type source \(D_R(M)\) and grade the coordinate ring \(C\) of its target by right translation. Pullback takes \(C_m\) into \(Rz^m\), with coefficient ideal \(I_m\subset R\). On a field fibre, a monic multiplicative-type homomorphism is closed: triangularize its diagonalizable action in a faithful representation; the occurring weights generate \(M\), since otherwise their common kernel would be nontrivial, and the matrix coefficients and the inverse determinant generate the full group algebra. Therefore the pullback on every field fibre is surjective. Weight projections commute with base change, so \(I_m+k(s)=k(s)\) at every prime. Locally some coefficient is consequently a unit, and \(I_m=R\). This holds for all \(m\), so \(C\to R[M]\) is surjective. Descent proves the criterion for a nonsplit source. Thus the lifted monomorphism is a closed immersion, with no unproved embedding criterion remaining.

To make the approximation step precise, use \textbf{Artin approximation for polynomial equations} in the following form. If \((R,\mathfrak m)\) is a Noetherian local G-ring, a solution \(\widehat a\in\widehat R^n\) of finitely many polynomial equations over \(R\) can, for every \(N\geq1\), be approximated modulo \(\mathfrak m^N\widehat R\) by a solution on an étale neighbourhood with the same residue field. Equivalently the approximating solution belongs to the henselization \(R^h\) and descends to a pointed étale neighbourhood. Comparison with the formal solution uses the canonical map from the chosen local étale algebra to \(R^h\to\widehat R\); no embedding of an entire, possibly disconnected, étale algebra is required.

A \textbf{G-ring} is a Noetherian ring whose local completion maps are regular, meaning flat with geometrically regular fibres. In particular the map \(R\to\widehat R\) above is regular.

Theorem 1.1 of the preceding \href{https://kokunoyumeto.github.io/open-math-courses/courses/AG-RG/AG-RG-S05.html}{Artin approximation lesson} proves this statement by its written desingularization sequence D.1--D.38, followed by the standard smooth residue-point lift in Lemma 1.3. The finite-order congruence is imposed in that same finite polynomial system. Its G-ring proof for every local ring essentially of finite type over \(\mathbb Z\) is F.132--F.135, including the completed-local faithful-flatness argument F.99A. These proofs retain arbitrary characteristic, imperfect residue fields, nilpotent elements and nonregular local G-rings.

Here is the finite system to which the theorem is applied. Choose finite presentations for the completed torus and a finite étale algebra splitting it. Since the completed base is local, that algebra is finite free; choose a basis and its multiplication table. Include the finitely many coefficients of the splitting isomorphism to a Laurent polynomial algebra, its inverse, the descent maps on the double overlap, and the homomorphism into \(G\). Each chosen algebra expression is a finite polynomial or Laurent polynomial. The algebra, Hopf, inverse-map and descent identities are tested on finite generating sets; equality in a finitely presented algebra is recorded using finitely many witnesses in its relation ideal. Invertibility conditions are encoded by adjoining inverse variables. In particular, invertibility of the trace-pairing determinant keeps the chosen finite free splitting algebra étale and faithfully flat.

The closed embedding must be kept in the same system. Its coordinate pullback is surjective by the preceding argument. Choose preimages of a finite set of generators of the completed torus algebra and include the equations stating that their images are those generators. After approximation, these equations still give a surjective pullback. The splitting maps and their inverses still identify the source with a split torus on the finite étale cover, and the preserved descent identities give a torus on the approximating base. Thus approximation retains both the torus structure and the closed embedding, rather than approximating only an arbitrary scheme map.

Apply the theorem at order one to this entire coefficient tuple. The solution then agrees with the specified fibre torus, its embedding and splitting data on the residue field; the equations themselves hold exactly on the étale neighbourhood. Spread its finite presentation and all chosen data from the local model to an open part of the Noetherian model, then pull back to the original base. This gives the required embedded torus after an étale neighbourhood, without a Noetherian assumption on \(S\). \(\square\)

\subsection{The positive-characteristic proof behind approximation}\label{AG-RG-01-the-positive-characteristic-proof-behind-approximation}

The full Néron--Popescu argument in \href{https://kokunoyumeto.github.io/open-math-courses/courses/AG-RG/AG-RG-S05.html}{Artin approximation}, Section 4, includes inseparable residue fields. The following elementary details correct its printed positive-characteristic construction and explain why that construction retains its full generality. We give them in our own words. The linked chapter is the Stacks Project authors' exposition, consulted in the AI Integrated Stacks fork; its \href{https://github.com/KokunoYumeto/unofficial-stacks-project-ai-drafts/blob/565b10e987aba5969b21145a0833f42d69f96790/README.md\#attribution-and-licensing}{attribution and GFDL 1.2 licence} remain those of the source.

Let \(k\) be a field of characteristic \(p>0\), and write \(k\to A\to\Lambda\), where \(A\) is of finite type over \(k\) and \(\Lambda\) is Noetherian and geometrically regular. Let \(H=H_{A/k}\) be the radical ideal defining the non-smooth locus, and choose a prime \(\mathfrak q\) minimal over \(\sqrt{H\Lambda}\). Put \(\Gamma=\Lambda_{\mathfrak q}\) and \(d=\dim\Gamma\). The construction chooses \(N\geq1\) with \(\mathfrak q^N\Gamma\subset H\Gamma\), then improves an auxiliary presentation and chooses an exponent \(c\geq1\). Its lifting step requires \(e=8c\). At the \href{https://github.com/KokunoYumeto/unofficial-stacks-project-ai-drafts/blob/565b10e987aba5969b21145a0833f42d69f96790/smoothing.tex\#L2687-L2730}{order-selection step}, take

\[
n=\max\{N+dc,\ d(e-1)+1\},\qquad e=8c.
\]

This replaces the printed \(n=N+dc\). The earlier choices of \(N\), the auxiliary algebras and \(c\) do not depend on \(n\). Define the finite set of coefficient images in \(\Gamma/\mathfrak q^n\Gamma\) with this order, then apply the \href{https://kokunoyumeto.github.io/open-math-courses/courses/AG-RG/AG-RG-S05.html}{Artinian approximation lemma}, D.35 with its global-selection proof D.34A, which permits every positive \(n\). All subsequent maps, Artinian subalgebras and parameters are chosen for this same quotient.

That lemma supplies a local polynomial ring \(P=k[y_1,\ldots,y_m]_{\mathfrak p}\) and a flat map \(P\to\Gamma\) for which \(\mathfrak p\Gamma=\mathfrak q\Gamma\). The chosen parameters \(\pi_1,\ldots,\pi_d\) generate \(\mathfrak pP\) and their images generate \(\mathfrak q\Gamma\). Put \(J=(\pi_1^e,\ldots,\pi_d^e)\subset P\). A degree-\(n\) monomial in \(d\) parameters contains some \(\pi_i^e\): otherwise every exponent is at most \(e-1\) and the total degree is at most \(d(e-1)<n\). Therefore

\[
\mathfrak p^nP\subset J,\qquad
\mathfrak q^n\Gamma\subset J\Gamma.
\]

This proves the \href{https://github.com/KokunoYumeto/unofficial-stacks-project-ai-drafts/blob/565b10e987aba5969b21145a0833f42d69f96790/smoothing.tex\#L3033-L3056}{inclusions used to descend the smooth factorization} from the quotient of order \(n\) to the quotient by \(J\). The unit rescalings of the parameters preserve these inclusions. In the auxiliary polynomial ring containing the \(t_i\), the selected target point sends \(t_i\) to a unit \(\delta_i\). Hence the ideals generated by \((\pi_i t_i)^e\) and by \(\pi_i^e\) agree after localization at that point. Reducing a smooth algebra modulo \(J\) is the required base change. The other lifting conditions involve \(e\), strict standardness and the annihilator equalities; these are unchanged by the larger order.

The coefficient calculation also survives that choice. At the \href{https://github.com/KokunoYumeto/unofficial-stacks-project-ai-drafts/blob/565b10e987aba5969b21145a0833f42d69f96790/smoothing.tex\#L2935-L2993}{denominator-clearing step}, write \(H=(a_1,\ldots,a_t)\) and \(\pi_i^N=\sum_j a_jd_j\) in the Artinian intermediate algebra \(D'\). Lift \(d_j\) to \(b_j\in\Gamma\). The error in this equality lies in \(\mathfrak q^n\Gamma\subset\mathfrak q^{n-N}H\Gamma\), so it is \(\sum_j a_jb'_j\) with \(b'_j\in\mathfrak q^{n-N}\Gamma\). Multiplication by \(\pi_i^N\) kills each \(b'_j\) modulo \(\mathfrak q^n\Gamma\). Choose a denominator unit \(s\in\Lambda\setminus\mathfrak q\); replace it by a suitable power and enlarge \(D'\) so that its residue also belongs to \(D'\), as explained below. Clearing denominators then gives coefficients \(\mu_{ij}\in\Lambda\) whose residues satisfy

\[
\mu_{ij}\bmod\mathfrak q^n\Gamma=s^{2N}\pi_i^Nd_j\in D'.
\]

The printed expression at line 2967 omits \(\pi_i^N\). This factor belongs to \(D'\) through the polynomial coordinate map, so the needed coefficient membership holds with the corrected expression. The subsequent unit and annihilator adjustments preserve the exact defining relations.

One further detail concerns \href{https://github.com/KokunoYumeto/unofficial-stacks-project-ai-drafts/blob/565b10e987aba5969b21145a0833f42d69f96790/smoothing.tex\#L2642-L2669}{enlarging the Artinian intermediate algebra}. Let \(\overline\Gamma\) be its Artinian target with maximal ideal \(\mathfrak n\) satisfying \(\mathfrak n^n=0\), and let \(\overline\lambda\) be a unit whose residue \(\alpha\) belongs to the residue field of \(D\). Choose a unit \(x\in D\) lifting \(\alpha^{-1}\) and put \(u=x\overline\lambda-1\in\mathfrak n\). For a power \(Q=p^r\geq n\), Frobenius gives

\[
(x\overline\lambda)^Q=(1+u)^Q=1+u^Q=1,
\qquad \overline\lambda^Q=x^{-Q}\in D.
\]

Thus the exponent must be a sufficiently large power of \(p\); divisibility by \(p\) alone does not give this identity. If \(\alpha\) is algebraic over the residue field \(F\) of \(D\), first take a \(p\)-power \(Q_0\) making \(\alpha^{Q_0}\) separable over \(F\). Lift that finite separable residue extension to the finite étale Artinian algebra \(D_1\) and its map to \(\overline\Gamma\). The same calculation with another sufficiently large \(p\)-power \(Q_1\) gives \(\overline\lambda^{Q_0Q_1}\in D_1\). The transcendental case uses the localized polynomial extension in the linked proof. These steps require no perfection assumption on \(k\) or on its residue extensions.

If \(d=0\), the localized regular target \(\Gamma\) is a field and the parameter lists are empty. Skip the sequence-lifting lemma, which is stated for a positive-length sequence, and use the Artinian smooth factorization and the height-zero delocalization step directly. The separable and inseparable cases therefore both cover dimension zero as well. These details are included in the preceding written desingularization proof; its Theorem D.38 and Theorem 1.1 supply the arbitrary-characteristic approximation input used above.

\textbf{Why fpqc tori split étale locally.} Suppose \(T/S\) is fpqc locally \(\mathbf G_m^r\). Affine descent makes it affine, smooth and of finite presentation, with geometric fibres tori; its relative dimension is locally constant. A fibre torus splits over a finite separable extension, by the separable-splitting proof in the earlier character lesson, with its full semilinear Hopf descent supplied in \href{https://kokunoyumeto.github.io/open-math-courses/courses/AG-RG/AG-RG-S08.html}{Supporting group-scheme proofs}, G.5.1--G.5.3. After making that residue extension, apply Lemma 3.1(3) with target \(G=T\) and source the entire fibre torus. We obtain an embedded étale torus \(T'\subset T\) near the point. Shrink so that both groups have the same relative dimension. On every geometric fibre the closed subtorus \(T'\) now has the dimension of the connected torus \(T\), so is the entire fibre. The inclusion between these smooth groups has invertible differential. \href{AG-RG-S02.html\#smooth-map-differential}{AG-RG-S02, Lemma 3.A}, makes it an étale monomorphism, hence an open immersion. Fibrewise surjectivity makes it an isomorphism. Finally split \(T'\) by its étale character cover. This proves the asserted equivalence of definitions, using the deformation lemma for a general smooth affine target, rather than presupposing the splitting assertion.

\subsection{Subgroups obtained by taking a limit}\label{AG-RG-01-subgroups-obtained-by-taking-a-limit}

Over a ring \(R\), let \(\lambda:\mathbf G_m\to H\) be a cocharacter of a smooth affine group. Define

\[
\begin{aligned}
P_H(\lambda)(A)&=\{h:\lambda(t)h\lambda(t)^{-1}\text{ extends to }t=0\},\\
U_H(\lambda)(A)&=\{h:\text{that extension has value }1\},\\
L_H(\lambda)&=C_H(\lambda(\mathbf G_m)).
\end{aligned}
\]

The extension is a morphism from \(\mathbf A^1_A\), not a numerical limit. All definitions commute with base change.

\textbf{Lemma 3.3.} These are smooth closed subgroups of finite presentation. The limit map identifies

\[
P_H(\lambda)=L_H(\lambda)\ltimes U_H(\lambda),
\]

and multiplication gives an open immersion

\[
U_H(-\lambda)\times P_H(\lambda)\longrightarrow H.
\]

If \(\mathfrak h=\bigoplus\mathfrak h_n\) is the cocharacter-weight decomposition, their Lie algebras are respectively \(\bigoplus_{n\ge0}\mathfrak h_n\), \(\bigoplus_{n>0}\mathfrak h_n\) and \(\mathfrak h_0\). The fibres of \(U_H(\lambda)\) are connected unipotent.

\textbf{Proof.} Give \(R[H]\) its conjugation grading, with \(f(t.h)=t^n f(h)\) for a homogeneous function of degree \(n\). Existence of the limit is equivalent to vanishing of all negative-degree functions at \(h\), so these functions generate the ideal of \(P_H(\lambda)\). Killing all nonzero-degree functions gives \(L_H(\lambda)\); adding to the ideal of \(P_H(\lambda)\) every degree-zero \(f-f(1)\) gives \(U_H(\lambda)\). These ideals are finitely generated over any \(R\). Indeed replace a finite algebra generating list of \(R[H]\) by all the finitely many homogeneous components appearing in it. A negative-weight monomial contains a negative-weight generator, so those generators generate the first ideal. The nonzero-weight generators generate the second. Modulo the first ideal, a weight-zero monomial involves only weight-zero generators: a product of remaining positive-weight generators cannot have weight zero. Thus the first ideal together with \(a-a(1)\) for the finitely many weight-zero generators \(a\) generates the third ideal. Adding these finite lists to a finite presentation gives finite presentation of all three subgroups. The coefficient projectors for the grading commute with every scalar extension, so these equations also prove arbitrary base-change compatibility.

Evaluation at zero is a homomorphism \(P\to L\) and is the identity on \(L\). The mutually inverse maps \((l,u)\mapsto lu\) and \(p\mapsto(\lim t.p,(\lim t.p)^{-1}p)\) prove the semidirect product assertion.

Here is a direct smoothness check over every base ring. The fixed subgroup \(L=H^{\mathbf G_m}\) is smooth by AG-GS-05, Theorem 3.4. To prove smoothness of \(P\), test a square-zero surjection \(C\to\bar C=C/J\). A point of \(P(\bar C)\) gives the equivariant extension \(F_0:\mathbf A^1_{\bar C}\to H_{\bar C}\) of its conjugation orbit. Smoothness of the affine \(H\) lifts it to a scheme map \(F:\mathbf A^1_C\to H_C\). On these maps let the universal \(a\in\mathbf G_m\) act by \[
(aF)(t)=\lambda(a)F(a^{-1}t)\lambda(a)^{-1}.
\] Since \(F_0\) is fixed, the ratio \((aF)F^{-1}\) is a regular one-cocycle in the vector module \[
V=\operatorname{Lie}(H_{\bar C})\otimes_{\bar C}J[t],
\qquad (av)(t)=\operatorname{Ad}(\lambda(a))v(a^{-1}t).
\] This square-zero kernel identification is made on the universal free algebras \(C[t,a,a^{-1}]\) and \(C[t,a,a^{-1},b,b^{-1}]\), so the displayed tensor ideal is exact there. The ratio identity \(c(ab)=c(a)+a c(b)\) follows from composition of this action; multiplication in the square-zero kernel is addition. AG-GS-05, Proposition 7.14 supplies \(w\in V\) with \(c(a)=w-aw\). Left multiplication of \(F\) by its kernel section \(w(t)\) therefore gives a fixed map, since \(aw+c(a)=w\). A fixed map is equivariant and, on \(\mathbf G_{m,C}\), is the conjugation orbit of its value at \(1\). Its extension to \(\mathbf A^1_C\) thus makes that value a lift in \(P(C)\).

For \(U\), start with \(F_0(0)=1\). Correct any scheme lift first by the inverse of its value at \(0\), a square-zero kernel element, so that \(F(0)=1\). Its equivariance defect now belongs to the stable submodule \(\operatorname{Lie}(H_{\bar C})\otimes_{\bar C}J[t]t\). Apply the same degree-one calculation to this submodule. The correcting \(w(t)\) vanishes at zero, so the resulting equivariant lift still has \(F(0)=1\) and gives a lift in \(U(C)\). We have proved formal smoothness of \(P\) and \(U\). Their already proved finite presentation and the infinitesimal criterion of AG-FSE, Infinitesimal lifting and invariance under thickenings, Theorem 4.1, give smoothness everywhere. Finally an identity tangent vector \(X=\sum_n X_n\) has orbit \(\sum_n t^nX_n\). It extends exactly when \(X_n=0\) for \(n<0\), and has value zero at the origin exactly when also \(X_0=0\). Fixedness means only weight zero remains. These are the three asserted Lie algebras.

The intersection \(U_H(-\lambda)\cap P_H(\lambda)\) is trivial: both positive and negative degrees must vanish, and the limit is the identity. The multiplication map is a monomorphism. Its differential is an isomorphism at the identity; at \((u,p)\), translations reduce it to the decomposition of \(\mathfrak h\) by \(\mathfrak u_-\oplus\mathfrak p\), followed by \(\operatorname{Ad}(p^{-1})\), which induces an automorphism modulo \(\mathfrak p\). \href{AG-RG-S02.html\#smooth-map-differential}{AG-RG-S02, Lemma 3.A}, therefore makes it étale everywhere; being a monomorphism, it is an open immersion.

Every point of \(U\) is connected to the identity by its extended orbit map \(\mathbf A^1\to U\), so the geometric fibres are connected. Over a field, embed \(H\rtimes\mathbf G_m\) faithfully in a general linear group and choose a weight basis for the acting torus. The condition that the limit be the identity makes \(U\) strictly triangular relative to that weight ordering. It is therefore unipotent. \(\square\)

\section{4. A torus over the ground field}\label{AG-RG-01-a-torus-over-the-ground-field}

The finite-field and infinite-field arguments have quite different mechanisms. The first adjusts descent by Lang's theorem. The second finds a semisimple infinitesimal direction and uses dimension induction. The infinitesimal argument is essential over imperfect fields.

\textbf{Lang's theorem in the form needed here.} If \(G/\mathbf F_q\) is smooth and connected and \(F\) is its Frobenius, the map \(g\mapsto g^{-1}F(g)\) is surjective on geometric points. By \href{AG-RG-S08.html\#field-group-components}{AG-RG-S08, Lemma G.0.3}, \(G_{\overline{\mathbf F}_q}\) is connected. Consider the action \(g\cdot x=gxF(g)^{-1}\) on that group. The differential of \(F\) is zero. Hence the orbit map at every point has invertible differential, and is étale by \href{AG-RG-S02.html\#smooth-map-differential}{AG-RG-S02, Lemma 3.A}. All its orbits are therefore open. Distinct orbits would give disjoint nonempty open subsets partitioning the connected variety, so there is only one. The orbit of the identity proves the assertion, after inversion. The same argument for the standard \(p\)-power map of \(\operatorname{GL}_n\) works over any algebraically closed field of characteristic \(p\).

\textbf{Jordan components and the height-one construction.} We give the algebraic tools in the infinite-field proof, so that no classification theorem is hidden in its infinitesimal step. Multiplicative Jordan components of a point of a closed linear algebraic group remain in that group. In characteristic \(p\), a large \(p\)-power kills its unipotent part; multiplication by that power is an automorphism on the reduced diagonalizable closure of the semisimple part's powers. The semisimple part therefore belongs to the closure of the powers of the original point, and so does its unipotent part. In characteristic zero, matrix entries of the powers are exponential polynomials \(\lambda^m P(m)\). Distinct nonzero \(\lambda\) give linearly independent such sequences, as repeated application of the difference operators \(f(m)\mapsto f(m+1)-\lambda f(m)\) shows. Consequently every polynomial equation holding on the powers holds separately on their diagonalizable and unipotent factors. This again puts both factors in their algebraic closure. The construction is compatible with homomorphisms, by uniqueness of the two matrix components.

Additive Jordan components of the Lie algebra of a smooth affine group also remain in its Lie algebra after algebraic closure. In characteristic \(p\), the Lie algebra is closed under \(p\)-th matrix powers. For a matrix \(X\), choose \(p^r\) large enough to kill its nilpotent part. A linearized polynomial recovers the semisimple part from \(X^{p^r}\). To verify this interpolation, let \(V\) be the \(\mathbf F_p\)-span of the eigenvalues raised to \(p^r\). The map \(\lambda^{p^r}\mapsto\lambda\) extends to an \(\mathbf F_p\)-linear map on \(V\), because Frobenius preserves exactly the \(\mathbf F_p\)-linear relations. If \(v_1,\ldots,v_d\) is a basis of \(V\), the equations \(\sum_{j=0}^{d-1}a_jv_i^{p^j}=\lambda_i\) have an invertible coefficient matrix: a nonzero linearized polynomial of degree at most \(p^{d-1}\) cannot vanish on the \(p^d\) elements of \(V\). Its solution gives the desired polynomial. Thus the semisimple part is a linear combination of higher restricted powers of \(X\), and its difference from \(X\) is the nilpotent part. In characteristic zero, the formal exponential of \(X\) lies in the group: a defining equation is annihilated by all powers of its invariant derivation and hence vanishes on the formal exponential. Its entries have the form \(e^{\lambda t}P(t)\). Applying differential operators \(d/dt-\lambda\) gives their independence for distinct \(\lambda\) just as for the preceding sequences. The algebraic closure of this formal one-parameter subgroup contains the separate semisimple and nilpotent one-parameter factors. Their derivatives are the required Lie components. Rationality over a perfect field follows from the polynomial construction and uniqueness. The same argument, or uniqueness in the adjoint representation, gives the adjoint Jordan components used below.

For a smooth group in characteristic \(p\), its first Frobenius kernel is encoded by the restricted Lie algebra \(\mathfrak g\). Explicitly, its distribution algebra is

\[
u(\mathfrak g)=T(\mathfrak g)/(xy-yx-[x,y],\ x^p-x^{[p]}).
\]

For an ordered basis, the restricted PBW monomials with exponents \(0,\ldots,p-1\) form a basis. Here is the argument in arbitrary characteristic. In the ordinary enveloping algebra, rewrite \(x_jx_i\) for \(j>i\) as \(x_ix_j+[x_j,x_i]\). Degree followed by the number of inversions decreases. The only overlapping reductions are triples of basis elements; their difference is the Jacobi identity, so ordered monomials give unique normal forms. Put \(z_i=x_i^p-x_i^{[p]}\). The identity \([x_i^p,y]=(\operatorname{ad}x_i)^p(y)=[x_i^{[p]},y]\) makes \(z_i\) central. Products of the \(z_i\) times ordered monomials with all exponents below \(p\) have distinct leading PBW monomials; division of each exponent by \(p\) proves that they form a basis of the ordinary enveloping algebra. Quotienting by all \(z_i\) therefore leaves exactly the asserted restricted basis. The restricted identities ensure that these basis relations impose \(x^p=x^{[p]}\) for every \(x\), not just the chosen basis. Its dimension is \(p^{\dim\mathfrak g}\). Local smooth coordinates at the identity show that the first Frobenius kernel has precisely this length; successive invariant differentiations pair these PBW monomials with its truncated coordinate monomials, with nonzero factorials of orders below \(p\). Hence the natural map of distribution algebras is an isomorphism. A restricted subalgebra \(\mathfrak h\subset\mathfrak g\) gives an injective Hopf map \(u(\mathfrak h)\hookrightarrow u(\mathfrak g)\) and, by duality, an actual closed height-one subgroup of \(G\). Thus both the positive-characteristic PBW assertion and its group correspondence have been supplied here.

If \(\mathfrak h\) is commutative and its semilinear \(p\)-operation is geometrically invertible, it has, over an algebraic closure, a basis fixed by that operation. Indeed write the operation as \(v\mapsto A v^{(p)}\) and apply the preceding Lang argument in \(\operatorname{GL}_d\) to change its matrix to the identity. In that basis, \(u(\mathfrak h)\) is the tensor product of the primitive Hopf algebras \(k[x]/(x^p-x)\); the dual groups are \(\mu_p\). Thus the subgroup is geometrically \(\mu_p^d\). Its Cartier dual is finite étale, so it is of multiplicative type over the original field and splits over a finite separable extension. This proves exactly the infinitesimal multiplicative-type assertion required in the argument.

\textbf{Lemma 4.1.} If a smooth connected affine \(k\)-group has a nontrivial \(k\)-torus, and geometrically maximal tori exist for groups of smaller dimension, then it has a geometrically maximal \(k\)-torus.

\textbf{Proof.} Let \(M\) be a nontrivial \(k\)-torus and put \(C=C_G(M)^0\). Smoothness is Lemma 3.1, and taking the identity component makes \(C\) geometrically connected by \href{AG-RG-S08.html\#field-group-components}{AG-RG-S08, Lemma G.0.3}. The quotient \(C/M\) is smooth connected affine and has smaller dimension. Affineness can be seen directly from the line detecting the normal subgroup \(M\): since \(M\) is central, it acts by the same scalar on the span of the translates of that line. The resulting representation in a projective linear group has kernel exactly \(M\), so its closed affine image is the quotient. By induction it has a geometrically maximal torus \(Q\). The inverse image \(E\) of \(Q\) is an extension of tori by tori, hence a torus: after an algebraic closure, the solvable structure theorem shows that its unipotent radical maps trivially to \(Q\) and lies in \(M\), so is trivial; it is a connected solvable reductive group, which is a torus.

Over an algebraic closure, a maximal torus of \(G\) containing \(M\) lies in \(C\). Thus the maximal torus dimensions in \(C\) and \(G\) agree. A torus of \(C\) properly containing \(E\) would yield a torus of \(C/M\) properly containing \(Q\), a contradiction. Therefore \(E\) is geometrically maximal in \(G\). \(\square\)

\textbf{Lemma 4.A (a commutative finite extension of multiplicative-type groups).} A finite commutative \(k\)-group \(E\) in an fppf exact sequence \[
1\longrightarrow N\longrightarrow E\longrightarrow Q\longrightarrow1
\] with \(N\) and \(Q\) of multiplicative type is of multiplicative type.

\textbf{Proof.} Extend to an algebraic closure; descent of multiplicative type is AG-GS-05, Theorem 5.0. For every \(E\)-representation \(W\), the \(N\)-fixed subspace is \(E\)-stable and the action factors through \(Q\). Factorization is an fppf descent assertion: the action is constant on the \(N\)-fibres, whose full torsor relation is supplied by the quotient theorem. It therefore descends through \(E\to Q\). Thus \[
W^E=(W^N)^Q.
\] Both fixed-subspace functors are exact by the earlier weight decomposition, so \(E\)-invariants are exact. Here is the full resulting splitting argument: for a subrepresentation \(V\subset W\), the surjection \(\operatorname{Hom}(W/V,W)\to\operatorname{Hom}(W/V,W/V)\) is equivariant. Exactness of invariants lifts the identity to an equivariant section, and induction decomposes every finite-dimensional representation into simples.

Write \(C=k[E]\) and \(H=C^\vee\). As \(E\) is finite commutative, \(H\) is a finite-dimensional commutative algebra; \(E\)-representations are precisely \(H\)-modules, by evaluating the coaction against this dual basis. The regular \(H\)-module is semisimple. A finite-dimensional commutative algebra with semisimple regular module is a product of fields: decompose into simple ideals; their annihilators are maximal, pairwise comaximal, and their intersection is zero, so the Chinese remainder map is an isomorphism. Over the algebraically closed field these fields are \(k\). The Hopf operations on \(H\) therefore describe a finite constant group \(M\), with \(H=k^M\). It is abelian because the comultiplication of \(H\) is dual to the commutative multiplication of \(C\) and is cocommutative. Dualizing the idempotent basis gives \(C=k[M]\), with its group-algebra Hopf operations. Hence \(E=D(M)\). \(\square\)

\textbf{Lemma 4.B (composite kernels stay central).} Suppose \(G\) is smooth geometrically connected, \(N\) is a finite connected central multiplicative-type subgroup, and \(Q\) is a finite connected central multiplicative-type subgroup of \(G/N\). Its inverse image \(E\) in \(G\) is finite connected central and of multiplicative type.

\textbf{Proof.} The quotient and its torsor map are the scheme quotient of AG-GS-04. Its restriction \(E\to Q\) is finite faithfully flat with kernel \(N\). Thus \(E\) is finite, and is connected: its geometrically connected fibres are torsors under the infinitesimal \(N\), and \(Q\) is geometrically connected.

Since \(Q\) is commutative and \(N\) is central, the commutator on \(E\) factors through a bilinear map \(Q\times Q\to N\). Independence of lifts and bilinearity follow from the elementary commutator identities when the commutators are central; the maps themselves descend through the fppf torsor covers. Such a bilinear map gives a morphism \[
Q\longrightarrow\underline{\operatorname{Hom}}(Q,N).
\] The target is the étale Hom scheme of AG-GS-05, Theorem 7.1, whose proof identifies it with the discrete character-homomorphism set after a separable splitting cover. A map from the connected infinitesimal \(Q\) to this étale scheme is constant, including on nilpotent test directions. Its value at the identity is zero. The commutator is therefore trivial and \(E\) is commutative. Lemma 4.A makes \(E\) of multiplicative type.

Conjugation now defines \(G\to\operatorname{Aut}(E)\), another étale scheme by the same Hom calculation. Geometric connectedness of \(G\) and the identity value force this morphism to be constant. Therefore \(E\) is central schematically, not merely on \(G(k)\).

Finally \(G/E\) is smooth geometrically connected affine. Affineness is AG-GS-04, Theorem 11.1d. After any algebraically closed scalar extension, its coordinate algebra injects into that of the smooth \(G\) by faithful flatness, and is reduced. AG-GS-02, Theorem 5.1 makes this reduced algebraic group smooth; AG-RG-S08, Lemma G.2.1 descends smoothness to \(k\). Its underlying space is the surjective image of the connected \(G\) after that extension. \(\square\)

\textbf{Lemma 4.C (descending the torus in a central affine group).} Let \(Z\) be an affine commutative algebraic \(k\)-group. The torus factor of \((Z_{\bar k}^0)_{\mathrm{red}}\) is defined over \(k\). In particular, a nontrivial such torus descends even if \(k\) is imperfect.

\textbf{Proof.} Let n run through all positive integers invertible in \(k\). The group \(Z[n]\) is finite étale. Indeed its tangent space at the identity is the kernel of multiplication by n on \(\operatorname{Lie}Z\), hence zero. Translation gives zero tangent space at every geometric point. A finite-type local algebra over an algebraically closed field with zero cotangent maximal ideal is a field by Nakayama. Thus \(Z[n]\) is geometrically reduced and zero-dimensional, hence finite étale; finite dimensionality and finite étaleness descend by the direct scalar-extension argument in AG-GS-05, Theorem 5.0.

Over a separable closure \(k_s\), these finite étale subgroups are unions of rational points. Let \(H\) be their union, and let \(Y_s\) be its reduced schematic closure in \(Z_{k_s}\). This closure is geometrically reduced and commutes with every further field extension. To see the exact reason, on an affine chart let J consist of the functions vanishing on \(H\). A scalar-extended function has a finite expression \[
\sum_i a_i f_i
\] with the \(a_i\) linearly independent over \(k_s\). It vanishes at every point of \(H\) precisely when each \(f_i\) does. Thus the evaluation kernel after an extension L is \(J\otimes L\). The quotient embeds by evaluation in a product of copies of L and is reduced. This argument establishes both the claimed kernel equality and geometric reducedness without assuming that \(k_s\) is perfect.

\(H\) is a subgroup. Its product is schematically dense in \(Y_s\times Y_s\) by the same coefficient-and-evaluation argument in two variables. Multiplication and inverse therefore restrict to \(Y_s\), since their equations vanish on \(H\times H\) and \(H\). The ideal J is Galois-stable. Its finite generating list and all its coefficients occur over a finite Galois extension; add their finitely many conjugates. Finite semilinear descent AG-RG-S08, Lemma G.5.2 descends the resulting ideal and its Hopf identities. Consequently \(Y_s=Y_{k_s}\) for a \(k\)-subgroup \(Y\). The base-change kernel equality just proved makes \(Y\) geometrically reduced.

Over \(\bar k\), the reduced connected central group is, by the already proved commutative solvable structure, \(T_Z\times U_Z\). An element of \(U_Z\) has no nontrivial prime-to-characteristic torsion: in positive characteristic its faithful unitriangular matrix has p-power order, and in characteristic zero the polynomial one-parameter subgroup from AG-RG-S03 Lemma 4.1 makes it torsion-free. Thus \[
H_{\bar k}\cap(Z_{\bar k}^0)_{\mathrm{red}}\subset T_Z.
\] All invertible-order torsion of \(T_Z\) lies in \(H\) and is schematically dense in \(T_Z\): a Laurent polynomial of finite support remains nonzero on \(T_Z[n]\) when n separates its exponent vectors. The identity component is an open and closed part of this finite-type group, so taking closure there commutes with its open restriction. It follows that \(Y_{\bar k}^0=T_Z\).

Here is the descent of this component explicitly. The reduced group \(Y_{\bar k}\) is smooth by AG-GS-02 Theorem 5.1; AG-RG-S08, Lemma G.2.1 descends smoothness to \(Y_{k_s}\). Smooth finite-type groups have finitely many disjoint irreducible components, so their identity component is open and closed. The extension \(\bar k/k_s\) is purely inseparable, and its projection is a homeomorphism. Here is the ring proof, including an infinite extension and nilpotents. In characteristic zero the extension is the identity. In characteristic \(p>0\), for any \(k_s\)-algebra \(A\), put \(B=A\otimes_{k_s}\bar k\). The inclusion \(A\to B\) is faithfully flat because extension of field scalars is exact and detects a nonzero vector space. Every \(b\in B\) is a finite sum \(\sum_i a_i\otimes\ell_i\). A common \(p\)-power sends all \(\ell_i\) into \(k_s\), so \(b^{p^r}\in A\). Thus \(B/A\) is integral. Lying over gives a prime above every prime \(\mathfrak p\subset A\). Such a prime is unique: membership of \(b\) in it is equivalent to membership of \(b^{p^r}\) in \(\mathfrak p\). For any ideal \(J\subset B\), apply lying over to \(A/(J\cap A)\subset B/J\): the image of \(V(J)\) is exactly \(V(J\cap A)\). Thus the map of spectra is closed. These lying-over and quotient arguments are proved in the earlier \href{AG-RG-S04.html}{AG-RG-S04 Lemma P0.4}. Hence this continuous bijection is a homeomorphism. The argument applies on every affine chart. Thus the identity component \(T_Z\) of \(Y_{\bar k}\) comes from an open and closed subscheme \(W_s\) of \(Y_{k_s}\), and its base change is exactly \(T_Z\) with its reduced scheme structure. The idempotent specifying \(W_s\) is fixed by every Galois automorphism, since the identity component is unique. Its finitely many coefficients lie in a finite Galois extension, so AG-RG-S08, Lemma G.5.2 descends that idempotent to \(k\). It defines an open and closed subscheme \(W\) of \(Y\). The unit and subgroup identities hold after the faithfully flat extension to \(\bar k\), so hold for \(W\); hence \(W\) is a \(k\)-subgroup with \(W_{\bar k}=T_Z\). AG-GS-05 Theorem 5.0 makes \(W\) separably split of multiplicative type, and its geometric character group is the free lattice of \(T_Z\). Therefore \(W\) is a \(k\)-torus. \(\square\)

\textbf{Theorem 4.2 (Grothendieck).} Every smooth connected affine group over a field \(k\) has a geometrically maximal torus defined over \(k\).

\textbf{Proof.} We use induction on dimension and give the imperfect-field step explicitly. The classical inputs about connected commutative groups and Lie algebras are specified at the end of the lesson.

If \(G_{\bar k}\) has a central maximal torus, it is unique, since all maximal tori are conjugate. Put \(Z=Z(G)\), with its scheme-theoretic structure. The torus factor of \((Z_{\bar k}^0)_{\mathrm{red}}\) is that maximal torus: it contains every central torus, and maximality leaves no larger torus. Lemma 4.C descends this factor to a \(k\)-torus. Its geometric fibre is the chosen maximal torus. This proves the theorem in the central-maximal-torus case, including the trivial-torus case; the same lemma will later descend a merely nontrivial central torus.

Suppose \(k\) is infinite and \(G_{\bar k}\) has a noncentral torus. We will find a nontrivial \(k\)-torus and then use Lemma 4.1. Choose a faithful representation and regard \(\mathfrak g=\operatorname{Lie}(G)\) inside a matrix algebra. Semisimple and nilpotent mean the additive Jordan components in this representation; these components belong to \(\mathfrak g_{\bar k}\) and are functorial in algebraic-group homomorphisms.

First assume that \(\mathfrak g(k)\) contains a semisimple element \(X\) with \(\operatorname{ad}(X)\ne0\). In characteristic zero its stabilizer \(C_G(X)\) is smooth by AG-RG-S08, Theorem G.3.4 applied to this finite-type subgroup scheme, and its Lie algebra is the kernel of \(\operatorname{ad}(X)\). It is proper and its identity component contains the nonzero semisimple tangent vector \(X\). After algebraic closure that component is not unipotent: Lemma 2.A makes its faithful Lie action strictly upper triangular if it were unipotent, which excludes \(X\). Dimension induction supplies a geometrically maximal \(k\)-torus in this smaller-dimensional smooth geometrically connected centralizer. Lemma 2.D makes its geometric fibre positive-dimensional, so the supplied \(k\)-torus is nontrivial.

In characteristic \(p>0\), a version of the same argument avoids smoothness questions for a Lie-element stabilizer. Put

\[
\mathfrak h=\operatorname{span}_k\{X,X^{[p]},X^{[p^2]},\ldots\}.
\]

This is a commutative restricted Lie subalgebra; its elements are commuting semisimple matrices, and its \(p\)-operation has no kernel after extending to \(\bar k\). The equivalence between restricted Lie algebras and finite height-one group schemes turns it into a subgroup \(H\subset\ker(F_{G/k})\) of multiplicative type. Indeed, over an algebraically closed field an invertible Frobenius-semilinear operation has a basis fixed by it, and each such line corresponds to \(\mu_p\), rather than \(\alpha_p\). Thus \(H_{\bar k}\simeq\mu_p^d\) with \(d>0\).

The centralizer \(C_G(H)\) is smooth by the invariant-cohomology argument of Lemma 3.1, which applies to all finite groups of multiplicative type. It is proper because the adjoint action on \(\mathfrak h\) is not trivial. Its identity component contains the connected \(H\). After algebraic closure that component cannot be unipotent, because the nontrivial multiplicative-type \(H\) cannot embed in a unitriangular group by Lemma 2.A. Dimension induction supplies a geometrically maximal \(k\)-torus in this smaller-dimensional smooth geometrically connected centralizer, and Lemma 2.D makes that torus nontrivial.

It remains to find a usable semisimple direction or to remove the central directions that obstruct this method. Choose a noncentral \(\mathbf G_m\subset G_{\bar k}\). Its adjoint action on \(\mathfrak g_{\bar k}\) has a nonzero weight, since otherwise its smooth centralizer would have dimension \(\dim G\). In characteristic zero its Lie generator has a nonzero eigenvalue under \(\operatorname{ad}\). Consequently the characteristic polynomial of \(\operatorname{ad}(Y)\) is not identically \(t^{\dim G}\) as \(Y\) varies in \(\mathfrak g\). Since \(k\) is infinite, some \(Y\in\mathfrak g(k)\) also has this property. Its semisimple component is rational over the perfect field \(k\) and has nonzero adjoint action. This is the situation just treated.

In characteristic \(p\), the same conclusion about a noncentral Lie generator can fail: all weights may have zero differential. If there is any semisimple rational vector with nonzero adjoint action, we are already done. Otherwise all semisimple rational vectors are central in \(\mathfrak g\). There is still a nonzero semisimple rational vector. A noncentral torus ensures that the Lie algebra is not a Lie algebra of nilpotent matrices. Hence some coefficient of the characteristic polynomial of the universal matrix \(Y\in\mathfrak g\) is nonzero. Choose a rational \(Y\) at which it is nonzero. Its semisimple part is nonzero and is defined over the perfect closure of \(k\). For sufficiently large \(r\), the matrix \(Y^{p^r}\) kills the nilpotent part and brings all coefficients of the semisimple part back into \(k\). It lies in \(\mathfrak g(k)\) by the restricted operation and is the required vector.

Choose one nonzero semisimple rational vector \(X\). It is central in \(\mathfrak g\) by the case assumption, and so are all its restricted powers, since \(\operatorname{ad}(X^{[p]})=(\operatorname{ad}X)^p\). Their span is a nonzero commutative restricted subalgebra with geometrically injective \(p\)-operation. Its height-one subgroup \(M\) is of multiplicative type. Its centralizer is smooth by exactness of multiplicative-type invariants, and has the whole Lie algebra \(\mathfrak g\): centralizing \(M\) is equivalent to acting trivially on its distribution algebra, generated by this central subalgebra. Thus \(C_G(M)\) has dimension \(\dim G\). Connectedness of \(G\) makes this centralizer all of \(G\), so \(M\) is central. No assertion about descent of the span of all semisimple vectors is needed.

Pass to \(G/M\) and repeat if all its rational semisimple tangent vectors are central. Lemma 4.B shows at each step that the composite kernel \(E_i\) is finite connected central and of multiplicative type, and that \(G/E_i\) is smooth geometrically connected affine. The kernel orders strictly increase, since each new height-one quotient kernel is nontrivial.

If the iteration continues indefinitely, its composite kernels \(E_i\) are central finite connected multiplicative-type groups by Lemma 4.B, with strictly increasing orders. Over \(\bar k\), write \(Z=Z(G)^0\) and \(R=Z_{\mathrm{red}}^0\). The quotient \(Z/R\) exists by AG-GS-04, Theorem 11.1d and is a finite group: the inclusion has the same underlying component and dimension, so the finite-type zero-dimensional quotient is finite. Let its order be c.~The restriction of \(Z\to Z/R\) to \(E_i\) has kernel \(E_i\cap R\). Its represented finite quotient maps monomorphically to \(Z/R\). A finite monomorphism is a closed immersion by AG-GS-03 Lemma 7.12's elementary proof: its field fibre is either empty or that field, and Nakayama kills the finite cokernel of its coordinate map. The quotient therefore has order at most c.~The finite torsor rank formula gives \[
|E_i|\le c\,|E_i\cap R|.
\] Thus some \(E_i\cap R\) is nontrivial. It is multiplicative type, as a subgroup of \(E_i\), and cannot lie in the unipotent factor of \(R\) by Lemma 2.A. The torus factor of \(R\) is therefore nontrivial. Lemma 4.C descends it.

If the iteration stops at a torus \(Q\) in \(G/E_i\), its inverse image \(P\) is a central extension of \(Q\) by infinitesimal \(E_i\). Its commutator descends to \(Q\times Q\to E_i\) and is zero because the smooth reduced source maps to the reduced identity of the infinitesimal target. Thus \(P\) is commutative. Its reduced identity group over \(\bar k\) surjects onto \(Q\): the finite infinitesimal torsor is a universal homeomorphism, and passing to reduction preserves its underlying image. Its unipotent factor maps trivially to \(Q\) by Lemma 2.A and so lies in the finite kernel; a positive-dimensional smooth connected unipotent group cannot do so. Its reduced identity group is consequently a nontrivial torus. Apply Lemma 4.C to the commutative affine \(k\)-group \(P\) itself. This descends that torus to a \(k\)-subgroup of \(P\subset G\). This finishes the infinite-field induction.

Finally let \(k=\mathbb F_q\), with Frobenius \(F\). Over \(\bar k\), choose a pair \((B,T)\) of a Borel subgroup and its maximal torus. Theorem 2.2 and Lemma 2.1 show that Borel--torus pairs are conjugate. Write \(F(B,T)=g(B,T)g^{-1}\). Lang's theorem gives \(h\in G(\bar k)\) with \(h^{-1}F(h)=g^{-1}\). Then

\[
F(h(B,T)h^{-1})=h(B,T)h^{-1}.
\]

Galois descent therefore gives a Borel subgroup and a geometrically maximal torus over \(k\). This closes the induction in every characteristic and over every field. \(\square\)

The infinitesimal step is needed because a noncentral group torus may have a central Lie algebra. In \(\operatorname{SL}_2\) in characteristic two, the Lie algebra of the diagonal torus consists of scalar matrices. The torus itself still acts nontrivially on the upper and lower root groups.

\section{5. The finite symmetry left after centralizing}\label{AG-RG-01-the-finite-symmetry-left-after-centralizing}

For a maximal torus in a reductive group scheme, define the fppf quotient

\[
W_G(T)=N_G(T)/T.
\]

The quotient has group structure because \(T\) is normal in its normalizer. It acts on \(T\), and thus on both character and cocharacter sheaves.

The general finite étale Weyl-scheme theorem is taught with its complete proof in \href{https://kokunoyumeto.github.io/open-math-courses/courses/AG-RG/AG-RG-04.html\#4-borel-subgroups-are-positive-systems}{Root data, Weyl chambers and the Bruhat decomposition, Theorem 4.2}, after the centralizer and reflection arguments it consumes. The matrix examples below have their own direct proofs here.

\textbf{Example 5.2.} For \(\operatorname{GL}_n\) and its diagonal torus, a normalizing matrix permutes the universal weight lines. Locally on the base it is a permutation matrix times an invertible diagonal matrix. Conversely every such matrix normalizes the torus. Hence

\[
N_{\operatorname{GL}_n}(D)=D\rtimes(\mathfrak S_n)_S,\qquad
W_{\operatorname{GL}_n}(D)=(\mathfrak S_n)_S.
\]

The permutation is locally constant; over a disconnected algebra it need not be a single global permutation.

\section{6. Exercises and solutions}\label{AG-RG-01-exercises-and-solutions}

\textbf{Exercise 6.1 (first steps).} Over an arbitrary scheme, compute the centralizer and normalizer of the diagonal torus of \(\operatorname{GL}_3\). Explain how the answer changes when the test algebra is disconnected.

\textbf{Solution.} The universal Laurent-polynomial computation in Example 1.1 makes the centralizer diagonal. A normalizer induces a permutation of the three rank-one character summands in the standard representation. After a decomposition by idempotents, the permutation is fixed on each piece. The normalizer is therefore \(D\rtimes(\mathfrak S_3)_S\), with quotient the constant finite étale group \(\mathfrak S_3\). Disconnected test algebras allow different permutations on different open and closed pieces, exactly as sections of a constant étale scheme do.

\textbf{Exercise 6.2 (descent).} Show that the two real tori in Example 1.2 are conjugate over \(\mathbb C\) but not over \(\mathbb R\). Show also that the multiplication torus \(\operatorname{Res}_{\mathbb C/\mathbb R}\mathbf G_m\) in \(\operatorname{GL}_{2,\mathbb R}\) is maximal and nonsplit.

\textbf{Solution.} The eigenvectors \((1,-i)\) and \((1,i)\) diagonalize the rotation matrix, with eigenvalues \(a+bi\) and \(a-bi\). On the norm-one subgroup these are inverse, so the eigenbasis carries it to the diagonal torus of \(\operatorname{SL}_2\). Rescale a column to obtain a determinant-one change of basis. Conjugation on its character lattice is \(m\mapsto-m\), whereas the diagonal real torus has trivial action, which rules out a real conjugation. For the multiplication torus, the two eigenvalues vary independently after complex base change; it becomes the full diagonal torus of \(\operatorname{GL}_2\). It is therefore maximal. Its two character generators are exchanged by conjugation, so it is nonsplit over \(\mathbb R\).

\textbf{Exercise 6.3 (finite fields).} For a smooth connected affine group over \(\mathbb F_q\), construct a Borel--torus pair over \(\mathbb F_q\) from one over \(\overline{\mathbb F}_q\). The only cohomological input permitted is Lang's theorem.

\textbf{Solution.} Frobenius carries the chosen pair to a conjugate pair, so choose \(g\) with \(F(B,T)=g(B,T)g^{-1}\). Solve \(h^{-1}F(h)=g^{-1}\) by Lang. Then \(F(hBh^{-1})=hBh^{-1}\) and \(F(hTh^{-1})=hTh^{-1}\). Their defining ideals descend by Galois descent. Their geometric fibres are the original Borel and maximal torus up to conjugation, so the descended pair has the required properties. Connectedness of \(G\) is a condition of Lang's theorem and cannot be removed from this argument.

\textbf{Exercise 6.4 (a family).} Over \(\mathbb Z\), let \(D\) be the diagonal torus in \(\operatorname{SL}_2\). Show that its nontrivial Weyl element is represented by \(\begin{pmatrix}0&-1\\1&0\end{pmatrix}\) even in characteristic two, and explain why the quotient is étale there.

\textbf{Solution.} Conjugation by the displayed matrix sends \(\operatorname{diag}(t,t^{-1})\) to \(\operatorname{diag}(t^{-1},t)\). A normalizer permutes the two universal weight lines, giving two components, each a coset of \(D\). Its quotient is the constant scheme with two points over \(\mathbb Z\), not \(\mu_2\). In characteristic two \(\mu_2\) is nonreduced, whereas a constant two-point scheme is still the disjoint union of two copies of the field and remains étale. The representative squares to \(-I\in D\); it need not itself have order two before taking the quotient.

\section{Exact supporting prerequisites and proof order}\label{AG-RG-01-exact-supporting-prerequisites-and-proof-order}

The classical triangularization, fixed-point, stabilizer, solvable structure and conjugacy arguments are in Section 2. Section 3 proves formal multiplicative-type homomorphism effectivity and the closed-immersion criterion. Section 4 proves Lang's theorem and supplies the Jordan and restricted-Lie arguments used in the field existence theorem. Their historical sources remain credited below.

The supporting geometric proofs are \href{https://kokunoyumeto.github.io/open-math-courses/courses/AG-GS/AG-GS-03.html}{Group schemes over a field}, for orbits and the closed-orbit lemma; \href{https://kokunoyumeto.github.io/open-math-courses/courses/AG-GS/AG-GS-04.html}{Quotients and torsors}, for finite free-action quotients; \href{https://kokunoyumeto.github.io/open-math-courses/courses/AG-RG/prerequisites.html\#prerequisite-ag-dfg-02}{Faithfully flat descent}, for affine schemes, representations and character descent; \href{https://kokunoyumeto.github.io/open-math-courses/courses/AG-GS/AG-GS-05.html}{Diagonalizable groups and groups of multiplicative type}, for the field character classification; \href{https://kokunoyumeto.github.io/open-math-courses/courses/AG-MO/AG-MO-03.html}{Limits of schemes and Noetherian approximation}, for finite-presentation spreading; the complete \href{https://kokunoyumeto.github.io/open-math-courses/courses/AG-RG/AG-RG-S05.html}{Artin approximation proof and application}, Sections 1--5, used in Section 3; and \emph{Dualizing sheaves and Serre duality for projective schemes}, for duality of curves. The prerequisite guide records their full statements and proof locations. The tori and reductive-group theorems assigned to this lesson are proved here.

The general cocharacter construction is proved here. The relative maximal-torus application is proved in the second lesson after the centralizer equality, and the finite Weyl-scheme theorem is proved in the fourth lesson after its reflection calculation. Their obligations are retained at those proof locations.

The \href{https://kokunoyumeto.github.io/open-math-courses/courses/AG-RG/prerequisites.html}{course prerequisite guide} records the exact supporting statements, their full proof routes, and the hypotheses needed in their applications.

\section{References}\label{AG-RG-01-references}

\begin{itemize}
\tightlist
\item
  Brian Conrad, \href{https://math.stanford.edu/~conrad/papers/luminysga3.pdf}{\emph{Reductive group schemes}}, in \emph{Autour des schémas en groupes}, Panoramas et Synthèses 42--43, Société Mathématique de France, 2014: §§2.1--2.3 and 3.2, Appendix A and §§B.2--B.3.
\item
  M. Demazure and A. Grothendieck, with the SGA 3 contributors, \href{https://webusers.imj-prg.fr/~patrick.polo/SGA3/}{\emph{Schémas en groupes}}, re-edited by P. Gille and P. Polo: Exposés III, VII A, IX--X, XII and XIV.
\item
  \href{https://www.jmilne.org/math/CourseNotes/iAG200.pdf}{Milne's freely accessible \emph{Algebraic Groups}, version 2.00 (2015)}. This exact free author edition provides comparison material; its citations do not replace the programme proofs.
\end{itemize}
\end{document}
